\begin{recipebox}{Mapo Tofu with Pork Mince}{4}{10 min prep + 15 min}

The Sichuan classic: soft tofu and pork in a spicy, savoury chilli-bean
sauce. The tofu does a lot of the protein work, so it uses less pork.
Roughly 37--40\,g protein per serve with rice.

\begin{ingredients}
  \item 300\,g pork mince
  \item 2 blocks (about 600\,g) firm or medium tofu, cut into 2\,cm cubes
  \item 2 tbsp chilli bean paste (doubanjiang)
  \item 3 garlic cloves, crushed
  \item 1 tbsp grated ginger
  \item 1 cup chicken stock
  \item 1 tbsp soy sauce
  \item 1 tsp sugar
  \item 1 tbsp cornflour mixed with 2 tbsp cold water
  \item 200\,g frozen peas or 2 heads bok choy
  \item \nicefrac{1}{2} tsp ground Sichuan pepper (optional)
  \item Spring onion, sliced
  \item 1\nicefrac{1}{4} cups dry rice, to serve
\end{ingredients}

\begin{method}
  \item Cook the rice.
  \item Fry the pork in a wok over high heat until browned and a little crispy.
  \item Add the chilli bean paste, garlic and ginger and fry for 1 minute until the oil turns red.
  \item Add the stock, soy and sugar and bring to a simmer. Gently slide in the tofu and peas and simmer for 4--5 minutes --- push the tofu around gently rather than stirring so it doesn't break up.
  \item Stir in the cornflour mixture and simmer for 1 minute until glossy and thick.
  \item Sprinkle with Sichuan pepper and spring onion and serve over rice.
\end{method}

\begin{notelist}
  \item Chilli bean paste (doubanjiang) is at Asian supermarkets. At a pinch, use 1 tbsp gochujang plus 1 tbsp soy sauce.
  \item Medium tofu is more traditional and silky; firm tofu holds together better.
\end{notelist}

\end{recipebox}
